% checklist.tex — ARR Responsible NLP Research Checklist

\section*{Responsible NLP Research Checklist}

\paragraph{A. Limitations.}
\begin{enumerate}[label=A\arabic*.,leftmargin=2em]
  \item Did you describe the limitations of your work? \textbf{Yes.} Section~\ref{sec:limitations} discusses dataset scale, lexical separability, annotation quality, OOD strategy prediction, single-seed variants, scope, and compute cost.
  \item Did you discuss any potential risks of your work? \textbf{Yes.} See below.
\end{enumerate}

\paragraph{B. Potential Risks and Ethical Considerations.}
This work studies multi-turn jailbreak attacks against LLMs. While the analysis is defensive (improving detection), we acknowledge dual-use risk: the persuasion taxonomy and attack strategy analysis could inform adversarial actors. We mitigate this by: (1) using only synthetic conversations generated by an open-weight model (Qwen3-8B), not real attack data targeting production systems; (2) not releasing attack-generation prompts that could enable novel attacks; (3) focusing the contribution on detection mechanisms rather than attack techniques. All jailbreak conversations are synthetic and do not contain real harmful content. No human subjects were involved.

\paragraph{C. Computational Experiments.}
\begin{enumerate}[label=C\arabic*.,leftmargin=2em]
  \item Did you report the number of parameters in the models used? \textbf{Yes.} DeBERTa-v3-base (184M) and RoBERTa-base (125M) in Sections~\ref{sec:method} and~\ref{sec:discussion}.
  \item Did you report the total amount of compute and the type of resources used? \textbf{Yes.} Under 4 GPU-hours total on a single NVIDIA A100 (40GB). Section~\ref{sec:limitations} and Appendix~\ref{sec:reproducibility}.
  \item Did you report the implementation details sufficient for reproducibility? \textbf{Yes.} Appendix~\ref{sec:reproducibility} provides all hyperparameters. Code and data are available at \url{https://anonymous.4open.science/r/persuasion_turn_jailbreak_detect-C207/}.
  \item Did you report error bars or confidence intervals? \textbf{Yes.} All OOD results report 3-seed mean$\pm$std (population standard deviation, $n$ divisor; Appendix~\ref{app:data_efficiency} uses Bessel-corrected $n{-}1$). Statistical tests (McNemar's, bootstrap CIs, permutation tests) are reported where applicable.
\end{enumerate}

\paragraph{D. Human Evaluation.}
Not applicable. No human evaluation was conducted. Strategy annotation quality was validated against a human sample ($n$=50, $\kappa$=0.78) but this was a validation check, not a full human evaluation study.

\paragraph{E. Existing Assets.}
\begin{enumerate}[label=E\arabic*.,leftmargin=2em]
  \item Did you cite the creators of artifacts you used? \textbf{Yes.} DeBERTa-v3~\citep{he2023debertav3}, RoBERTa~\citep{liu2019roberta}, Qwen3-8B~\citep{qwen3_2025} (for data generation).
  \item Did you discuss the license or terms for use and/or distribution of any artifacts? \textbf{Yes.} The training and evaluation data are synthetically generated and will be released under the CC BY-NC 4.0 license. Code will be released under MIT. DeBERTa-v3 is available under MIT; RoBERTa under MIT; Qwen3-8B under Apache 2.0.
  \item Did you discuss whether the data you are using contains personally identifiable information? \textbf{N/A.} All data is synthetically generated. No PII is present.
\end{enumerate}
